\documentclass[10pt,a4paper]{article}

\usepackage[utf8]{inputenc}
\usepackage[T1]{fontenc}
\usepackage[scaled=0.98]{helvet}
\renewcommand{\familydefault}{\sfdefault}
\usepackage[top=0.82cm,bottom=0.82cm,left=1.12cm,right=1.12cm]{geometry}
\usepackage{xcolor}
\usepackage{hyperref}
\usepackage{enumitem}
\usepackage{microtype}
\usepackage{titlesec}

\input{glyphtounicode}
\pdfgentounicode=1

\pagestyle{empty}
\setlength{\parindent}{0pt}
\setlength{\parskip}{0pt}
\setlength{\emergencystretch}{2em}
\urlstyle{same}

\definecolor{ink}{HTML}{161616}
\definecolor{accent}{HTML}{17365D}
\definecolor{muted}{HTML}{555555}

\hypersetup{
  colorlinks=true,
  urlcolor=accent,
  linkcolor=accent,
  pdfauthor={Lucas Gonzalez Fiz},
  pdftitle={Lucas Gonzalez Fiz - LLM AI Engineer CV}
}

\titleformat{\section}
  {\large\bfseries\color{accent}}
  {}{0pt}{}
  [\vspace{-4pt}\color{accent}\rule{\linewidth}{0.45pt}]
\titlespacing*{\section}{0pt}{5.8pt}{3pt}

\setlist[itemize]{
  leftmargin=1.15em,
  itemsep=1.1pt,
  topsep=1.5pt,
  parsep=0pt,
  partopsep=0pt
}

\newcommand{\entry}[4]{
  \textbf{#1}\hfill{\small #2}\\[-1pt]
  \textit{#3}\hfill{\small\textit{#4}}\par\vspace{-1.5pt}
}

\newcommand{\project}[3]{
  \textbf{#1} \textbar\ {\small\color{muted}\textit{#2}}\\[-1.5pt]
  #3\par\vspace{2pt}
}

\begin{document}

{\fontsize{26.5pt}{28pt}\selectfont\bfseries\color{ink} Lucas Gonz\'alez Fiz}\par
\vspace{1pt}
{\fontsize{12.5pt}{14pt}\selectfont\bfseries\color{accent} AI Engineer | LLM, RAG \& Agentic Systems}\par
\vspace{3pt}
{\small
Ourense, Spain \textbar\ +34 677 228 205 \textbar\
\href{mailto:lucasglezfiz04@gmail.com}{lucasglezfiz04@gmail.com}\\[-1pt]
\href{https://linkedin.com/in/lucas-gonzalez-fiz}{linkedin.com/in/lucas-gonzalez-fiz} \textbar\
\href{https://github.com/LucachuTW}{github.com/LucachuTW} \textbar\
\href{https://lucasfiz.dev}{lucasfiz.dev}
}

\section{Summary}
AI Engineer focused on LLM applications, RAG, GraphRAG and governed agentic systems. Experienced in access-controlled retrieval, hybrid search, reranking, tool authorization, human approval workflows, evaluation gates and FastAPI-based deployment, with additional applied-ML experience in healthcare and computer vision.

\section{Education}
\entry{B.Sc. in Artificial Intelligence (240 ECTS)}{Graduated Jul 2026}{University of Vigo -- ESEI}{Ourense, Spain}
\textbf{Selected coursework:} Natural Language Processing, Information Retrieval, Knowledge Representation, Semantic Web, Machine Learning I--II, Big Data and Distributed Systems.\\[-1pt]
\textbf{Bachelor's thesis (12 ECTS):} Designed a reproducible multi-seed benchmark comparing PPO, A2C, DQN and RecurrentPPO in partially observable environments, with systematic experiment configuration and evaluation.\\[-1pt]
\textbf{Academic focus:} Retrieval, language technologies, knowledge-based systems, machine learning evaluation and responsible AI.\\[-1pt]
\section{Activities \& Awards}
\textbf{Formula Student AI:} Silverstone competitor \textbar\ \textbf{HackUDC 2026:} Inditex Tech Challenge podium finisher.

\section{Skills}
\textbf{LLM \& Retrieval:} RAG, GraphRAG, LangGraph, MCP, Hugging Face, Qdrant, Neo4j, BM25, reranking\\
\textbf{Agents \& Backend:} tool calling, approval workflows, authorization, auditability, FastAPI, PostgreSQL\\
\textbf{ML \& Data:} Python, SQL, PyTorch, scikit-learn, embeddings, evaluation pipelines\\
\textbf{Infrastructure:} Docker, Kubernetes, Linux, Git, Ollama, local model serving

\section{Experience}
\entry{AI Software Developer -- Autonomous Systems}{Nov 2024 - Present}{Auria Formula Student AI Team}{Ourense, Spain}
\begin{itemize}
  \item Develop real-time perception software on a 1/10-scale autonomous racing platform for the Auria Formula Student AI team, working in an international Scrum environment.
  \item Train and benchmark pretrained and scratch YOLO cone detectors, applying structured pruning and TensorRT-oriented export for embedded deployment.
  \item Integrate camera and perception components into reproducible Linux, Docker and Git workflows, with ROS/NATS services prepared for NVIDIA Jetson deployment.
\end{itemize}

\entry{AI Engineer Intern}{Mar 2026 - Jun 2026}{MicroPort CRM}{Ourense, Spain}
{\small\color{muted}\textit{Contracted through Cardiovascular Gallega}}\par\vspace{-1.5pt}
\begin{itemize}
  \item Built an LLM/RAG pipeline to normalize heterogeneous clinical tables and map terminology across medical data sources, combining retrieval, prompting and validation stages.
  \item Developed and evaluated a machine-learning model for detecting pacemaker-lead malfunctions to support clinical review workflows.
\end{itemize}

\section{Selected Projects}
\project{\href{https://lucasfiz.dev/en/projects/airlock-agents}{Airlock Agents}}{LangGraph, MCP, PostgreSQL, Ollama, FastAPI, Kubernetes}{Built a governed multi-agent platform with role-scoped tools, read-only SQL access, durable human approval gates and end-to-end audit traces; added offline regression gates over documents and \textbf{541,909} retail transactions.}

\project{\href{https://lucasfiz.dev/en/projects/strata}{Strata}}{GraphRAG, Qdrant, Neo4j, BGE-M3, Qwen3, LangGraph}{Applied clearance filters inside vector and graph retrieval before generation; combined dense and BM25 search, graph expansion and reranking, improving exact-ID recall from \textbf{0.79 to 1.00}.}

\project{\href{https://lucasfiz.dev/en/projects/sentinel-vision}{Sentinel Vision}}{PyTorch, FastAPI, Docker, TensorRT}{Delivered an observable real-time inference service with explicit performance evaluation, API deployment and hardware-aware optimization, demonstrating end-to-end ML engineering beyond LLM applications.}

\section{Languages}
Spanish -- Native \textbar\ English -- B2 CEFR

\end{document}
